\begin{abstract}
\fablenote{AAAI repositioning: AAAI is a broad AI venue, not an NLP venue. Lead with the decision-theoretic story (cost-sensitive information acquisition / VOI under uncertainty) and present the linguistic finding as the surprising enabler, not the headline. Suggested opening: ``Acquiring information has a cost; acting on too little of it has another. We study when an AI system should pay for higher-fidelity perception...'' This frames VOILA as resource-bounded reasoning, a classic AAAI theme, while keeping the linguistic insight as the key technical surprise.}
Multimodal retrieval-augmented systems store images at multiple fidelities---captions, thumbnails, compressed JPEGs, full-resolution originals---yet retrieve at fixed resolution regardless of the query, paying maximum acquisition cost even when a coarser representation would suffice.
We formalize \emph{pre-retrieval fidelity selection}: given only the question text, select the cheapest fidelity whose expected accuracy justifies its retrieval cost, \emph{before any image is accessed}.
This problem is distinct from post-retrieval cascades (FrugalGPT, RouteLLM) that act after context is fetched, and from concurrent resolution-selection methods that still require image access.
We show the problem is tractable from question text alone: linguistic structure (e.g., \emph{how many}, \emph{what color} vs.\ \emph{is there}, \emph{does it show}) strongly predicts required visual fidelity.
\textsc{VOILA} estimates per-fidelity success probabilities via a calibrated gradient-boosted classifier over question features, then applies a greedy value-of-information rule---with no image access and no VLM execution at decision time.
Across five datasets and six VLMs (7B--235B parameters), \textsc{VOILA} reduces retrieval cost by 50--65\%, matching or exceeding full-resolution on OCR-heavy and domain-specialized benchmarks while exposing a controllable accuracy--cost trade-off on general VQA benchmarks.
\fablenote{The final sentence ends the abstract on the paper's weakest number (``losing at most 15 accuracy points''). Reframe the trade-off as a \emph{tunable operating point}: VOILA exposes $\lambda$ as an accuracy--cost dial, so the honest and stronger claim is ``traces the accuracy--cost Pareto frontier: at matched accuracy it cuts cost by X\%, and at matched cost it gains Y points over the strongest router baseline.'' Also fix the inconsistency: this abstract says 50--65\% but other parts of the project say 50--60\%; pick one number set and use it everywhere (abstract, intro, conclusion).}

% -----------------------------------------------------------------------
% OLD ABSTRACT (kept for reference)
% -----------------------------------------------------------------------
%Despite significant costs from retrieving and processing high-fidelity visual inputs, most multimodal vision-language systems operate at fixed fidelity levels. We introduce VOILA, a framework for Value-Of-Information-driven adaptive fidelity selection in Visual Question Answering (VQA) that optimizes \emph{what information to retrieve} before model execution. Given a query, VOILA uses a two-stage pipeline: a gradient-boosted regressor estimates correctness likelihood at each fidelity from question features alone, then an isotonic calibrator refines these probabilities for reliable decision-making. The system selects the minimum-cost fidelity maximizing expected utility given predicted accuracy and retrieval costs. We evaluate VOILA across three deployment scenarios using five datasets (VQA-v2, GQA, TextVQA, LoCoMo, FloodNet) and six Vision-Language Models (VLMs) with 7B--235B parameters. VOILA consistently achieves 50--60\% cost reductions while retaining 90--95\% of full-resolution accuracy across diverse query types and model architectures, demonstrating that pre-retrieval fidelity selection is vital to optimize multimodal inference under resource constraints.

\end{abstract}
